\documentclass[acmsmall,nonacm,screen]{acmart}

\usepackage{booktabs}
\usepackage{array}

\citestyle{acmauthoryear}
\settopmatter{printacmref=false}

\AtEndPreamble{\theoremstyle{acmdefinition}\newtheorem{remark}[theorem]{Remark}}

\newcommand{\kw}[1]{\mathsf{#1}}
\newcommand{\str}[1]{\texttt{#1}}
\newcommand{\far}{d}
\newcommand{\Down}{\kw{Down}}
\newcommand{\Level}{\kw{Level}}
\newcommand{\Finish}{\kw{Finish}}

% A sentence that states a fact in words. The script has one check for every name.
\newcommand{\claim}[2]{#2}

% The rows of a table that the script writes.
\makeatletter
\newcommand{\rows}[1]{\@@input artifact/expected/#1.tex }
\makeatother

% written by budget.py; do not edit
\newcommand{\exAutomata}{5,000}
\newcommand{\exBudgetLoose}{40}
\newcommand{\exBudgetTight}{eight}
\newcommand{\exConfigurations}{16,184}
\newcommand{\exCutLooseRuns}{7}
\newcommand{\exCutTight}{100}
\newcommand{\exCutTightRuns}{2,000}
\newcommand{\exDead}{1,284}
\newcommand{\exExampleCharacters}{22}
\newcommand{\exExampleFar}{eight}
\newcommand{\exExampleStates}{30}
\newcommand{\exExampleTokens}{33}
\newcommand{\exExampleWidest}{nine}
\newcommand{\exForbidden}{29,676}
\newcommand{\exLowerFails}{77}
\newcommand{\exMutants}{seven}
\newcommand{\exOutputs}{22,413}
\newcommand{\exPricedOne}{10,655}
\newcommand{\exPricedTwo}{5,903}
\newcommand{\exPushdown}{600}
\newcommand{\exRuns}{2,000}
\newcommand{\exSafe}{3,080}
\newcommand{\exSize}{800}
\newcommand{\exSpanning}{1,970}
\newcommand{\exTouchedLooseRuns}{13}
\newcommand{\exUpperLoses}{23}
\newcommand{\exUsualFails}{1,198}
\newcommand{\exWithLetters}{1,813}


\begin{document}

\title{A Budget Is a Distance: Structured Output That Finishes in Time}

\author{Claude-Fable-5.2}
\authornote{An AI model. The proofs in this paper are on paper and have not been checked by a
machine. A script compares the decoder with what trying every sequence of tokens gives, on
small automata and vocabularies. At the time of writing nobody has reviewed the paper.}

\begin{abstract}
Constrained decoding makes the output of a language model follow a schema. At every step it
forbids the tokens after which no valid output could follow. It does not make the output
finish. A model has a limit on the tokens that it may emit, and when the limit is reached in
the middle of a string or an array, what arrives is a valid beginning and nothing more. We
close this gap with one number for every state of the schema's automaton: its distance, the
fewest tokens of the vocabulary that complete the output from there. A decoder that allows a
token only if the distance after it is no more than the tokens that remain always ends with a
complete output, whatever the model chooses, and it loses nothing: every sequence of tokens
within the limit that spells a valid output is still possible. No mask that guarantees
completion allows more. The distance of the start state is the least limit that a request
with the schema can be given, known before a token is generated. The distance must be counted
in tokens. Counted in characters, it gives a decoder that either fails or forbids outputs
without need. The same holds for a budget in any other unit, such as characters, and for
two budgets at once. For regular schemas the distance is a shortest path, and the budget
changes only the last few tokens of an output. For nested schemas a token can close several levels at
once, and the distance of a configuration is computed from summaries by one pass over its
stack. A script checks all of this against brute force on \exAutomata{} random automata and
\exPushdown{} random pushdown schemas, and it notices \exMutants{} deliberate mistakes.
\end{abstract}

\keywords{constrained decoding, structured output, grammars, shortest paths, pushdown systems, tokens, language models}

\authorsaddresses{}

\maketitle

\section{Introduction}

A program that calls a language model usually wants an answer that it can parse: an object
with certain fields, a call to a tool with arguments of certain types. Constrained decoding
gives it one. The schema is compiled to an automaton, and at every step of the generation
the decoder forbids the tokens that the automaton cannot read, or after which it could never
reach an accepting state. Whatever the model writes is then a prefix of a valid output, and
when the model stops, the output is valid
\citep{Willard2023,Geng2023,Koo2024}.

The last clause hides a condition. The model stops by its own choice only where the schema
allows it. But every request also has a limit on the number of tokens, and when that limit
is reached, the generation is cut where it stands. If that is in the middle of a string, the
caller receives a prefix of a valid output, which is not a valid output and does not parse.
The documentation of one vendor says so in as many words: if the answer is cut at the limit,
it may be incomplete and may not match the schema, and the remedy is to ask again with a
higher limit \citep{claudestructured}. For a caller this is the worst of both: the tokens are
paid for, and the guarantee for which constrained decoding was chosen does not hold.

The gap has a simple reason. The decoder asks of every token whether the output \emph{can}
still be completed after it. With a limit, the question is whether it can be completed
\emph{in the tokens that remain}. That is a question of distance. For every state of the
automaton, let its distance be the fewest tokens that lead from it to an accepting state. A
decoder that knows the distances allows a token only if the distance after it is at most the
number of tokens left after it. Far from the limit this forbids nothing more than before.
Near the limit it steers the model towards an end, and at the last token it allows only what
closes the output.

This paper works the idea out and proves what it promises.

\begin{itemize}
\item \emph{The budget mask} (Section~\ref{sec:mask}). If the limit is at least the distance
  of the start state, every run of the decoder ends with a complete output, whatever the
  model chooses (Theorem~\ref{thm:safe}). Nothing is lost: every sequence of tokens within
  the limit that spells a valid output is a run (Theorem~\ref{thm:lost}). No mask that
  guarantees completion allows a token that this one forbids (Theorem~\ref{thm:greatest}). A
  limit below the distance of the start state cannot be met, and the request can be refused
  before it costs anything (Theorem~\ref{thm:before}). For regular schemas the budget changes
  the mask only in the last $D$ steps, where $D$ is the largest distance of a state
  (Theorem~\ref{thm:end}).
\item \emph{Tokens, not characters} (Section~\ref{sec:characters}). The distance has to be
  measured in the tokens of the vocabulary. The distance in characters is easier to compute
  and gives two tempting masks, of which one loses outputs and the other fails
  (Proposition~\ref{prop:characters}).
\item \emph{Nested schemas} (Section~\ref{sec:nested}). For schemas with a stack the
  distance depends on the whole stack, and a token such as \str{]\}} closes two levels at
  once, so the distance is not a sum over the levels. It is computed from summaries of the
  schema, with one step for every entry of the stack (Theorem~\ref{thm:nested}). The mask
  then bounds the depth of nesting by what the budget can still close
  (Corollary~\ref{cor:depth}).
\item \emph{Other budgets} (Section~\ref{sec:other}). A limit on characters, or on any price
  that tokens have, is met in the same way with the distance in that unit
  (Theorem~\ref{thm:priced}). Two limits at once need more than one number per state, and a
  small table is enough (Theorem~\ref{thm:two}).
\end{itemize}

Section~\ref{sec:example} shows an example, Section~\ref{sec:experiments} the tests, and
Section~\ref{sec:limits} the limits. The most important one is that the mask guarantees the
form of the output and not its sense. An output that was steered to an end is complete, and
a list in it may be shorter than the model wanted. A harness should say when that happened,
and the decoder knows.

\section{An Example}\label{sec:example}

A schema asks for an object with a name and a list of tags. To keep the example small, names
and tags are words over the letters \str{a} and \str{b}:
\begin{verbatim}
    {"name":"ab","tags":["a","ba"]}
\end{verbatim}
The automaton of this schema has \exExampleStates{} states. The vocabulary has
\exExampleTokens{} tokens: every character that the schema uses, and pieces such as
\str{\{"}, \str{name}, \str{":"}, \str{","}, \str{tags}, \str{":[}, \str{"]\}} and \str{ab},
as a tokenizer trained on such text would have them.

\claim{example}{The shortest valid output, with a name of one letter and no tags, has
\exExampleCharacters{} characters. In tokens it is much shorter: \exExampleFar{} tokens
spell it, and no fewer do. So the distance of the start state is \exExampleFar{}. No state
is further than \exExampleWidest{} tokens from an accepting state.}

Now let the limit be nine tokens, and let the model write as Table~\ref{tab:trace} shows.
After eight tokens it has opened the first tag. Its ninth token is another letter of that
tag, and with it the limit is reached. The caller receives \str{\{"name":"ab","tags":["ab},
which no parser accepts.

\begin{table}
\caption{A run of nine tokens that the usual mask allows. After each token: the tokens that
are left, the distance of the state in tokens and in characters, and whether the budget mask
allows the token, which it does if the distance is at most what is left.}\label{tab:trace}
\begin{tabular}{rlrrrl}
\toprule
 & token & left & distance & in characters & allowed \\
\midrule
\rows{trace}
\bottomrule
\end{tabular}
\end{table}

\claim{example-trace}{The budget mask allows the first eight tokens and forbids the ninth:
after it the distance would be one with no token left. With the token \str{"]\}} in its
place, the output is complete.} \claim{example-cut}{At that step the usual mask allows
thirteen tokens. With one token left the budget mask allows only \str{"]\}}, and with two
tokens left it allows eleven, among them every letter: the tag may still grow by one token,
if the token after that closes everything.} So the mask does not end the output early. It
ends it at the last moment at which it can still be ended.

The column of distances in characters shows why the count must be in tokens. After the
eighth token three characters are missing and one token. A decoder that counted characters
would have had to stop the model long before.

\section{The Budget Mask}\label{sec:mask}

\subsection{Schemas, Tokens and Decoders}

A \emph{schema} is a deterministic automaton $(Q, \Sigma, \delta, q_0, F)$ over an alphabet
$\Sigma$ of characters, with a partial transition function $\delta$ and a set $F$ of
accepting states. Its language is the set of valid outputs. A \emph{vocabulary} $V$ is a
finite set of non-empty strings over $\Sigma$, the \emph{tokens}. We write $\delta(q, t)$ for
the state that the characters of a token $t$ lead to from $q$, if all the transitions exist.
The \emph{token graph} has the states as its nodes and an edge from $q$ to $\delta(q, t)$
for every token $t$ for which that state exists.

The \emph{distance} $\far(q)$ of a state is the length of a shortest path in the token graph
from $q$ to an accepting state, and $\infty$ if there is none. So $\far(q) = 0$ exactly for
accepting states, and $\far(q)$ is the fewest tokens that complete an output from $q$.

A \emph{decoder} is at a \emph{point} $(q, r)$: the state $q$ that the output so far leads
to, and the number $r$ of tokens that it may still emit. It starts at $(q_0, B)$, where $B$
is the \emph{budget}. A \emph{mask} $M$ gives a set of tokens $M(q, r)$ for every point with
$r \ge 1$. At a point the model either stops, which it may do only if $q$ is accepting, or
picks a token $t$ of $M(q, r)$, and the decoder goes to $(\delta(q, t), r - 1)$. We assume
nothing about how the model picks. A \emph{run} is a sequence of points that arises in this
way. The decoder \emph{fails} at a point $(q, r)$ if $q$ is not accepting and it can go
nowhere: $r = 0$, which is an output cut at the limit, or $M(q, r)$ is empty. A mask is
\emph{safe} for a budget if no run from $(q_0, B)$ reaches a point where the decoder fails.

The \emph{usual mask} allows the tokens after which an output can still be completed:
$M(q, r) = \{t : \delta(q, t) \text{ can reach an accepting state}\}$, whatever $r$ is. The
\emph{budget mask} is
\[
  M_B(q, r) \;=\; \{\, t \in V \;:\; \far(\delta(q, t)) \le r - 1 \,\}.
\]

\subsection{What the Budget Mask Guarantees}

\begin{theorem}[Safety]\label{thm:safe}
If $\far(q_0) \le B$, the budget mask is safe for $B$: every run ends in an accepting state
after at most $B$ tokens, or can be continued.
\end{theorem}

\begin{proof}
Every point $(q, r)$ of a run has $\far(q) \le r$. This holds at the start by assumption, and
the mask keeps it. Let $(q, r)$ be a point with $q$ not accepting. Then $\far(q) \ge 1$, so $r
\ge 1$. The first token $t$ of a shortest path from $q$ leads to a state at distance $\far(q)
- 1 \le r - 1$, so $t$ is in $M_B(q, r)$. The decoder does not fail.
\end{proof}

\begin{theorem}[Nothing is lost]\label{thm:lost}
Let $t_1, \ldots, t_n$ be tokens with $n \le B$ that spell a valid output. Then the decoder
with the budget mask can emit them and stop. So its outputs are exactly the sequences of at
most $B$ tokens that spell a valid output.
\end{theorem}

\begin{proof}
Let $q_i$ be the state after the first $i$ tokens. The tokens $t_{i+1}, \ldots, t_n$ lead
from $q_i$ to an accepting state, so $\far(q_i) \le n - i \le B - i$. The decoder is at
$(q_{i-1}, B - i + 1)$ when it comes to $t_i$, and the mask there asks that $\far(q_i) \le B
- i$. After $t_n$ the state is accepting, and the decoder may stop. Conversely the decoder
stops only in an accepting state and after at most $B$ tokens.
\end{proof}

The two theorems together say that the budget mask conditions the decoder on one event,
that the output is complete within the budget, and on nothing else. The next theorem says
that no mask can do better.

\begin{theorem}[The greatest safe mask]\label{thm:greatest}
Let $M$ be a mask that is safe for $B$. If $M$ allows a token at a point that one of its runs
reaches, the budget mask allows that token there.
\end{theorem}

\begin{proof}
Let $M$ allow $t$ at $(q, r)$, and suppose that $\far(\delta(q, t)) > r - 1$. Continue the
run from $(\delta(q, t), r - 1)$ with tokens that $M$ allows, never stopping, for as long as
possible. This ends after at most $r - 1$ tokens, at a point with no tokens left or with an
empty mask. Its state is not accepting, because no path of at most $r - 1$ tokens leads from
$\delta(q, t)$ to an accepting state. So the decoder fails there, and $M$ is not safe.
\end{proof}

\begin{theorem}[Before the first token]\label{thm:before}
If $\far(q_0) > B$, no sequence of at most $B$ tokens spells a valid output, and no mask is
safe for $B$.
\end{theorem}

\begin{proof}
The first claim is the definition of the distance. A safe mask would have a run, and the run
would end in an accepting state after at most $B$ tokens.
\end{proof}

So $\far(q_0)$ is the least limit that makes sense for a schema and a vocabulary. A service
can compute it when it compiles the schema and reject a request with a lower limit at once,
with a message that names the number. Today such a request is run, paid for, and returns an
incomplete output.

\begin{theorem}[Only near the end]\label{thm:end}
Let $D$ be the largest distance of a state that has one. At a point with $r > D$, the budget
mask allows exactly the tokens after which some sequence of tokens completes the output.
\end{theorem}

\begin{proof}
If $\far(\delta(q, t))$ is finite, it is at most $D \le r - 1$.
\end{proof}

A decoder with the budget mask therefore behaves like the usual one until at most $D$ tokens
are left, and the budget can change no more than the last $D$ tokens of an output. In the
example $D$ is \exExampleWidest{}.

\subsection{Cost}

The distances are found by one search that starts from the accepting states and follows the
edges of the token graph backwards. The edges are what a constrained decoder computes in any
case, when it works out for every state which tokens it allows \citep{Willard2023}. The
search adds time linear in the number of edges, which is at most $|Q| \cdot |V|$, and one
number for every state. For the mask, the tokens that a state allows are grouped by the
distance of the state that they lead to. The mask at $(q, r)$ is the union of the groups
with distance at most $r - 1$, and a state has at most $D + 1$ groups. With more than $D$
tokens left the mask is the one that the decoder has already.

\section{Tokens, Not Characters}\label{sec:characters}

The distance of a state in characters, $c(q)$, is the length of a shortest string that leads
from $q$ to an accepting state. It needs no vocabulary and is the first thing that comes to
mind. It is the wrong number.

\begin{proposition}[Characters]\label{prop:characters}
Let $\ell$ be the length of the longest token.
\begin{enumerate}
\item $c(q) \le \ell \cdot \far(q)$ for every state. If every character is a token, then
  $\far(q) \le c(q)$, and a state can be completed by tokens exactly if it can be completed
  by characters.
\item If some character is not a token, a state can be completed by characters and not by
  tokens, and the usual mask, which judges by characters, can then fail whatever the budget
  is.
\item Let every character be a token. The mask $\{t : c(\delta(q, t)) \le r - 1\}$ is safe
  for every budget of at least $c(q_0)$, and it can forbid tokens that the budget mask
  allows. The mask $\{t : \lceil c(\delta(q, t)) / \ell \rceil \le r - 1\}$ allows every
  token that the budget mask allows, and it is not safe.
\end{enumerate}
\end{proposition}

\begin{proof}
(1) A path of $\far(q)$ tokens spells at most $\ell \cdot \far(q)$ characters, and if every
character is a token, a string is a path of as many tokens. (2) and the negative halves of
(3) are shown by examples, below and in Section~\ref{sec:experiments}. For the first mask of
(3), the proof of Theorem~\ref{thm:safe} applies with $c$ in place of $\far$, because a
character is a token; and $c \ge \far$ makes the mask a part of the budget mask. For the
second, $\lceil c / \ell \rceil \le \far$ by (1).
\end{proof}

\paragraph{A vocabulary that cannot spell.} \claim{stuck}{Let the schema accept only
\str{ab}, and let the tokens be \str{a}, \str{ab} and \str{bb}. After the token \str{a} the
output can be completed by a character and by no token. The usual mask allows \str{a}, and
then it allows nothing, so the decoder fails. The budget mask allows only \str{ab} at the
start.} Real vocabularies contain every byte, so every string can be spelled. The example
matters for vocabularies that have been restricted, and it shows that the right notion of
``can be completed'' is a path in the token graph even without a budget.

\paragraph{The mask by characters.} In the example of Section~\ref{sec:example} the start
state is \exExampleCharacters{} characters and \exExampleFar{} tokens from the end.
\claim{example-characters}{With a limit of nine tokens, the first mask of part~3 allows
nothing at the start: after the best first token twenty characters are still missing, and
the mask takes each of them for a token. The second mask divides by the longest token, which
has four characters, and takes the start state to be six tokens from the end. With a limit of
six or seven it lets the model begin, and no output of so few tokens exists.}

\section{Nested Schemas}\label{sec:nested}

A schema for nested data is not regular. An array can contain arrays, and the automaton needs
a stack to know how many brackets are open.

\subsection{Pushdown Schemas}

A \emph{pushdown schema} has states $Q$ with a start state and accepting states, an alphabet
$\Sigma$, and a finite set $\Gamma$ of stack symbols. It reads one character at a time. A
move depends on the state, the character and the symbol on top of the stack, or the fact
that the stack is empty, which we write as the top $\bot$. It leads to a state and keeps the
stack, pushes a symbol, or pops the top. A \emph{configuration} $(q, \gamma)$ is a state and
a stack, and it is accepting if $q$ is and $\gamma$ is empty. Nested data fits this form: an
opening bracket pushes, the closing bracket pops, and the state says, for instance, whether
the output is inside a string, where a bracket is a character like any other.

A token moves a configuration by all its characters. The token graph now has the
configurations as its nodes, infinitely many, and the distance $\far(q, \gamma)$ is again
the length of a shortest path to an accepting configuration. The decoder, its points, the
budget mask and safety are defined as before, with configurations in place of states. The
proofs of Theorems~\ref{thm:safe} to~\ref{thm:before} use nothing about the token graph but
shortest paths, so the four theorems hold for pushdown schemas as they stand.
Theorem~\ref{thm:end} does not: the deeper the stack, the further the end, and there is no
largest distance.

\subsection{The Distance of a Configuration}

What remains is to compute $\far(q, \gamma)$. One would like to add up, for every entry of
the stack, the tokens that it takes to close that level. That is wrong, because a token can
close several levels: \str{]\}} closes an array and the object around it, and \str{"]\}}
ends a string as well. The price of closing a level depends on how much of a token is
already in hand when the level is reached.

A \emph{point} is a pair $(q, w)$ of a state and a proper prefix $w$ of a token: a token has
been begun and $w$ of it has been read. Between two tokens $w$ is empty. Reading a character
$a$ at the point $(q, w)$, with a move of the schema from $q$ to $q'$, leads to the point
$(q', wa)$ at price $0$ if $wa$ is a proper prefix of a token, and to the point $(q',
\varepsilon)$ at price $1$ if $wa$ is a token. Both can be the case. The price of a path of
such steps is the number of tokens that it completes.

For a top $f$, a symbol or $\bot$, let $\Level_f(x, y)$ be the least price of a path from
the point $x$ to the point $y$ that begins and ends with $f$ on top of the same stack and
never makes that stack shorter. These numbers are the least solution of three rules. A point
reaches itself at price $0$. A path can be extended by a move that keeps the stack. And a
path can be extended by a move that pushes a symbol $g$, a path of $\Level_g$, and a move
that pops $g$. For a symbol $f$, let $\Down_f(x, y)$ be the least price of a path of
$\Level_f$ from $x$ followed by a move that pops $f$ and ends at $y$. Let $\Finish(x)$ be
the least $\Level_\bot(x, (q, \varepsilon))$ over the accepting states $q$.

\begin{theorem}[Nested distance]\label{thm:nested}
For a configuration with the state $q$ and the stack $f_1 \ldots f_k$, with $f_1$ on top,
\[
  \far(q, f_1 \ldots f_k) \;=\; \min_{y_1, \ldots, y_k}\;
  \Down_{f_1}((q, \varepsilon), y_1) + \Down_{f_2}(y_1, y_2) + \cdots + \Down_{f_k}(y_{k-1}, y_k)
  + \Finish(y_k),
\]
where the $y_i$ range over the points.
\end{theorem}

\begin{proof}
A path of tokens from the configuration to an accepting one is a path of steps on points
that begins and ends between two tokens, and its length in tokens is its price. Cut it after
the first move that makes the stack shorter than $k$ symbols, after the first that makes it
shorter than $k - 1$, and so on. The $i$-th piece begins with $f_i$ on top, never makes the
stack shorter until its last move, and ends with the move that pops $f_i$: its price is at
least $\Down_{f_i}$ between its two ends. The last piece begins and ends with an empty stack
in an accepting state between two tokens, so its price is at least $\Finish$. Conversely,
paths that realize the terms of the sum fit together to a path of tokens. That $\Level$ is
given by the three rules is the usual decomposition of a path that returns to its level into
moves on the level and excursions above it, as in the summaries of interprocedural analysis
and of weighted pushdown systems \citep{Reps2005}.
\end{proof}

The sum is evaluated from the bottom of the stack. With every entry of the stack the decoder
keeps a vector: for every point, the least price of finishing everything from that entry
down. A push computes one new vector from the one below, a pop drops one, and the distance
of the current configuration is one entry of the top vector. The tables $\Level$ and $\Down$
are computed once for a schema and a vocabulary, by iterating the three rules.

\begin{corollary}[Depth]\label{cor:depth}
At every point that a run of the budget mask reaches, with $k$ symbols on the stack and $r$
tokens left, $k \le \ell \cdot r$, where $\ell$ is the length of the longest token.
\end{corollary}

\begin{proof}
A token pops at most $\ell$ symbols, so $\far(q, \gamma) \ge k / \ell$, and the points of a
run have a distance of at most $r$.
\end{proof}

So the budget also limits how deep the model may nest, and it does so by itself: an opening
bracket is forbidden when the brackets that would then be open could no longer all be
closed.

\subsection{An Example: Lists of Lists}

Let a value be the letter \str{a} or a list of values in square brackets, separated by
commas, as in \str{[[a,a],[]]}. The schema has three states, for the places where a value
must come, where a list has just been opened, and where a value has ended, and one stack
symbol for an open bracket. The vocabulary has the four characters and the tokens \str{[[},
\str{]]}, \str{],[}, \str{a]} and \str{a,a}. Table~\ref{tab:lists} gives the distance of a
configuration by the number of brackets that are open.

\begin{table}
\caption{Lists of lists: the distance in tokens by the number of open brackets, in the three
states of the schema. A list cannot have just been opened with no bracket
open.}\label{tab:lists}
\begin{tabular}{rrrr}
\toprule
open brackets & after a value & a list just opened & before a value \\
\midrule
\rows{lists}
\bottomrule
\end{tabular}
\end{table}

\claim{lists}{After a value, $k$ open brackets are closed with $\lceil k/2 \rceil$ tokens,
because \str{]]} closes two of them, and the same holds where a list has just been opened.
Where a value must come, the distance is $1 + \lfloor k/2 \rfloor$: the token \str{a]} writes
the value and closes a bracket. A sum over the levels would give $k$ and $k + 1$. With a
budget of two, four and six tokens, the decoder with the budget mask opens at most two, four
and six brackets.}

\section{Other Budgets}\label{sec:other}

\subsection{A Budget in Another Unit}

Limits are not always counted in tokens. A text message has a number of characters, a column
of a database a number of bytes, the first line of a commit message by convention fifty characters,
and a bill is counted in money. Let every token $t$ have a \emph{price} $w(t) > 0$ in some
unit, and let the budget $B$ be an amount in that unit. Let $\far_w(q)$ be the least total
price of a path of tokens from $q$ to an accepting state. The decoder is at a point $(q, r)$
with the amount $r$ left, and a token $t$ takes it to $(\delta(q, t), r - w(t))$.

\begin{theorem}[Priced budgets]\label{thm:priced}
The mask $\{t : w(t) + \far_w(\delta(q, t)) \le r\}$ is safe for $B$ if $\far_w(q_0) \le B$.
Its outputs are exactly the sequences of tokens of total price at most $B$ that spell a valid
output, and no safe mask allows a token that it forbids. If $\far_w(q_0) > B$, no mask is
safe.
\end{theorem}

\begin{proof}
The proofs of Theorems~\ref{thm:safe} to~\ref{thm:before} go through with $r - w(t)$ in place
of $r - 1$ and paths of least price in place of shortest paths. Runs are finite because
prices are positive.
\end{proof}

The distances $\far_w$ are found by the algorithm of \citet{Dijkstra1959} on the token graph
with its edges reversed. With the length of a token as its price, $\far_w$ is the distance
in characters, counted along tokens. So Section~\ref{sec:characters} has a second reading:
the distance in characters was wrong there because the budget was in tokens. \emph{The
distance must be measured in the unit of the budget.}

\subsection{Two Budgets at Once}

A request can have both limits, at most $B$ tokens and at most $P$ in price. One number per
state is then not enough, because a state can have a completion that is short and dear and
another that is long and cheap. Let $f_n(q)$ be the least total price of a path of at most
$n$ tokens from $q$ to an accepting state:
\[
  f_0(q) = 0 \text{ if $q$ is accepting, and } \infty \text{ otherwise;} \qquad
  f_n(q) = \min\bigl(f_{n-1}(q),\ \min_{t}\, w(t) + f_{n-1}(\delta(q, t))\bigr).
\]
The decoder is at a point $(q, r, p)$ with $r$ tokens and the amount $p$ left.

\begin{theorem}[Two budgets]\label{thm:two}
The mask $\{t : f_{r-1}(\delta(q, t)) \le p - w(t)\}$ is safe for the budgets $B$ and $P$ if
$f_B(q_0) \le P$. Its outputs are exactly the sequences of at most $B$ tokens and of total
price at most $P$ that spell a valid output, and no safe mask allows a token that it
forbids. If $f_B(q_0) > P$, no mask is safe.
\end{theorem}

\begin{proof}
Every point of a run has $f_r(q) \le p$. If $q$ is not accepting, then $r \ge 1$, and the
first token $t$ of a cheapest path of at most $r$ tokens has $w(t) + f_{r-1}(\delta(q, t)) =
f_r(q) \le p$, so the mask allows it. A sequence within both budgets is a run, because after
its first $i$ tokens the rest of it is a path of at most $B - i$ tokens whose price is at
most what is left. A token that the mask forbids has no completion within what is left of
both budgets, and a mask that allows it fails as in Theorem~\ref{thm:greatest}.
\end{proof}

The table $f$ has one row for every number of tokens up to $B$, and its rows stop changing
after at most $|Q|$ steps. A decoder that looked only at the cheapest completion of a state,
whatever its length, would be wrong, and the script has that mistake among the ones that it
plants.

\subsection{In a Service}

Three uses of the distance need no change to the model.

\emph{Refuse early.} By Theorem~\ref{thm:before} a request whose limit is below
$\far(q_0)$ cannot succeed. The service knows $\far(q_0)$ when it has compiled the schema,
and it can answer with an error that names the least limit.

\emph{Reserve.} A response often has free text, or reasoning, before its structured part,
under one limit. Seen as one schema, the free part is a state that loops and from which the
structured part starts, so its distance is that of the structured part, or one more. The budget mask then
ends the free part when what is left is just enough for the shortest structured output, and
nothing has to be reserved by hand.

\emph{Report.} When the budget removes a token that the usual mask allows, the output is no
longer what the model would have written. The decoder sees the first step at which this
happens and can return it with the response, as a reason for stopping of its own, next to
``the model stopped'' and ``the limit was reached''.

\section{Experiments}\label{sec:experiments}

The script \texttt{budget.py} implements the decoder for regular and for pushdown schemas.
It trusts the functions that run an automaton on a string and a pushdown schema on a token,
and it compares everything else with what trying every sequence of tokens gives. It needs
nothing but Python, and \texttt{make -C artifact check} runs it and compares its output with
the numbers in this paper.

\paragraph{Regular schemas.} The script draws \exAutomata{} random automata with two to six
states over two or three letters, each with a vocabulary of a few tokens of up to three
letters, half of the time with every letter among them. \claim{regular}{For every state the
distance equals the fewest tokens found by trying every sequence of up to four. For a random
budget of up to four tokens, the distance of the start state says correctly whether any
output fits. Where one fits, in \exSafe{} cases, no run under the budget mask fails. The
outputs under the budget mask are exactly the accepted sequences within the budget,
\exOutputs{} in all. Each of the \exForbidden{} tokens that the mask forbids at some point
has no completion in what is left, and with more tokens left than the largest distance the
mask allows every token that has a completion at all. There are \exDead{} states that can be
completed by characters and not by tokens, and the usual mask has a failing run in
\exUsualFails{} of the cases in which an output fits.}

\paragraph{Characters.} \claim{characters}{In \exWithLetters{} of the cases every character
is a token and an output fits the budget. The mask from the distance in characters loses
outputs in \exUpperLoses{} of them, and the mask from that distance divided by the longest
token has a failing run in \exLowerFails{}.} The first number is small because random
automata with a few states rarely have a long way to the end; the example of
Section~\ref{sec:example} is more typical of schemas with keywords.

\paragraph{Other budgets.} For every automaton the script also draws a price of one to three
for every token and a budget of up to five in that unit. \claim{priced}{The decoder of
Theorem~\ref{thm:priced} has exactly the accepted sequences within that budget as its
outputs, \exPricedOne{} in all, and the decoder of Theorem~\ref{thm:two}, with the budget in
tokens as well, exactly those within both, \exPricedTwo{} in all. Neither has a run that
fails where an output fits, and each says before the first token whether one fits.}

\paragraph{Nested schemas.} The script draws \exPushdown{} random pushdown schemas with two
or three states, one or two stack symbols, and two to four tokens of one or two characters.
\claim{nested}{For each of the \exConfigurations{} configurations with up to three symbols
on the stack, the distance from the summaries equals the result of a search over
configurations with up to five tokens; \exSpanning{} configurations with a stack that is not
empty have a completion. Under the budget mask with a budget of up to four tokens, no run
fails where an output fits, the outputs are the accepted sequences within the budget, and
the stack is never deeper than Corollary~\ref{cor:depth} allows.}

\paragraph{A model that picks at random.} How often the limit cuts an output depends on the
model and on the task, and we have no real model here. To see the mechanism at work, the
script runs the example of Section~\ref{sec:example} with a model that picks, with equal
chances, among everything that it may do next, including to stop where that is allowed.
Table~\ref{tab:runs} compares the two masks for six budgets, \exRuns{} runs each.

\begin{table}
\caption{The example with a model that picks at random. For each budget: the percentage of
runs under the usual mask that end incomplete, the same under the budget mask, the
percentage of runs under the budget mask in which the mask differs from the usual one at some
step, the percentage of steps at which it differs, and the most tokens that are left at
such a step.}\label{tab:runs}
\begin{tabular}{r@{\qquad}rr@{\qquad}rrr}
\toprule
 & \multicolumn{2}{c}{incomplete} & \multicolumn{3}{c}{the budget mask differs} \\
budget & usual & budget & runs & steps & tokens left \\
\midrule
\rows{runs}
\bottomrule
\end{tabular}
\end{table}

\claim{runs}{With the least possible budget of \exBudgetTight{} tokens, \exCutTightRuns{} of
\exRuns{} runs under the usual mask end incomplete. With \exBudgetLoose{} tokens,
\exCutLooseRuns{} still do. Under the budget mask no run ends incomplete at any budget. The
two masks never differ with more than \exExampleWidest{} tokens left, as
Theorem~\ref{thm:end} says, and with \exBudgetLoose{} tokens the budget mask differs from the
usual one in \exTouchedLooseRuns{} runs.} A model that picks at random writes longer lists
than a real model would, so the shares on the left are no forecast. The right half shows
what Theorem~\ref{thm:end} promises: with a generous budget the mask is almost never seen.

\paragraph{Planted mistakes.} A check that never fails checks nothing. The script has
\exMutants{} switches, each of which plants one mistake: a mask that allows a token when the
distance after it equals the tokens left before it; a distance measured from the start
state; a token that moves the automaton by its first character only; a decoder that may stop
wherever the output can still be completed; a decoder with two budgets that looks only at
the cheapest completion; a nested distance in which no token reaches
across a pop; and a pass over the stack from the wrong end. \claim{all-noticed}{Each of them
is noticed by a comparison that fails.}

\section{Related Work}\label{sec:related}

\paragraph{Constrained decoding.} \citet{Willard2023} compile a regular expression or a
grammar to an automaton and index, for every state, the tokens that it allows, so that the
mask costs a lookup. \citet{Koo2024} recast this with automata theory and treat the
tokenizer as a transducer. \citet{Geng2023} argue for grammars as a general interface to
structured tasks. PICARD rejects tokens with an incremental parser \citep{Scholak2021},
Synchromesh derives the allowed continuations from a completion engine
\citep{Poesia2022}, and SynCode, DOMINO and XGrammar make grammar masks fast and precise for
full programming languages and for pushdown structure
\citep{Ugare2024,BeurerKellner2024,Dong2024}. All of these guarantee that the output so far
can be completed. As far as we know, none of them relates the completion to the tokens that
remain, and the vendor documentation cited in the introduction describes the consequence.
The construction of this paper adds one number per state to any of these systems.

\paragraph{What masking does to the distribution.} A mask changes the distribution of the
model, and renormalizing at every step is not the same as conditioning on a valid output
\citep{Park2024}. The budget mask is a mask and shares the problem: near the limit it
shifts weight to short completions. Theorem~\ref{thm:end} bounds where this happens, and
Theorems~\ref{thm:lost} and~\ref{thm:greatest} say that the set of possible outputs is as
large as it can be. How to sample from that set in proportion to the model is a separate
question.

\paragraph{Shortest completions.} The cheapest way to complete a prefix of a sentence is an
old tool of syntax-error repair. \citet{Roehrich1980} builds parsers that can always
continue a valid prefix to a sentence, and \citet{BurkeFisher1987} choose among repairs by
cost. Our distance is the length of the cheapest continuation, with two differences: it is
counted in the tokens of a vocabulary, which cut across the terminals of the grammar, and it
is used before the fact, to keep the generator within reach of the end.

\paragraph{Shortest paths and pushdown systems.} For regular schemas the distance is a
shortest path in a graph with unit weights \citep{Dijkstra1959}. For pushdown schemas,
Section~\ref{sec:nested} is an instance of weighted pushdown systems \citep{Reps2005}, with
tokens completed as the weight. What is specific here is the product with the begun token,
which is what lets a token reach across a pop.

\paragraph{Budgets.} An earlier paper of this collection bounds the tokens that an agent
loop spends with sized types \citep{tpp6}, and another one types programs that generate and
then check \citep{tpp3}. This paper is about a single generation and makes the check
unnecessary for one kind of failure.

\section{Limits}\label{sec:limits}

\paragraph{Form, not sense.} The mask guarantees an output that parses. It does not
guarantee that the output says what the model was going to say. A list can end after three
entries because the budget ended, and a string can stop in the middle of a word. This is
still better than an output that does not parse, but the caller must be told. The decoder
knows the first step at which the budget removed a token that the usual mask allowed, and it
should report that, as a reason for stopping of its own.

\paragraph{The distribution.} As said in Section~\ref{sec:related}, the mask distorts the
model's distribution in the last steps. We prove what is possible, not how likely it is.

\paragraph{Schemas.} We treat regular schemas and pushdown schemas that read one character
per move. Conditions that are not of this kind, such as that the entries of a list differ,
are outside. A grammar with empty moves or with lookahead has to be brought into this form
first.

\paragraph{Real vocabularies and real models.} The tests use vocabularies of a few tokens
and a model that picks at random. We argue that the distance costs little for a vocabulary
of a hundred thousand tokens, because it reuses the edges that every constrained decoder
computes, but we have not measured it. For pushdown schemas the tables grow with the number
of begun tokens that matter at a pop, and we have not measured that either. We have no data
on how often real outputs are cut at the limit.

\paragraph{What counts as a token.} We count the tokens of the output and let the model
stop for free. A service that counts a stop token, or that shares one limit between thinking
and output, has to shift the budget accordingly.

\paragraph{Proofs on paper.} The proofs are short, and none has been checked by a machine.
The script tests their statements on small instances.

\section{Conclusion}

Constrained decoding promises a valid output and delivers a valid prefix. The difference is
a distance, the fewest tokens to the end, and it is cheap to know: one number per state for a
regular schema, one vector per stack entry for a nested one. With it, a limit on tokens stops
being a way for a structured output to fail. The decoder steers the last tokens towards an
end, the service can name the least limit for a schema, and a request below it is refused
before it is paid for.

\bibliographystyle{ACM-Reference-Format}
\bibliography{refs}

\end{document}
